\subsection{Exact algebraic spectral defects of the source relation}
\label{sec:cq-sourcealgebra}
These calculations keep the actual $\cqN$ and require only the
already proved $M\cqN\subseteq\cqN$. For every $\lambda\in\cqC$,
division by $x=s-\lambda$ gives the algebraic identities
\begin{equation}\label{eq:cq-source-kercoker}
 \begin{aligned}
 \operatorname{coker}_{\rm alg}(S-\lambda)
     &=\cqE/(\cqN+x\cqE),\\
 \ker(S-\lambda)&\simeq
        (\cqN\cap x\cqE)/(x\cqN),
        & [n]&\longmapsto[n/x].
 \end{aligned}
\end{equation}
Division belongs to the original order class: if $f(\lambda)=0$,
the filled quotient is entire, is bounded on a fixed disk, and
outside $|s|\geq2|\lambda|+1$ its modulus is at most $2|f(s)|$.
Thus $x\cqE$ is exactly the kernel of evaluation at $\lambda$.
The cokernel identity follows from the image of the induced
multiplication. For the second identity, changing $n$ by $xn_0$
changes $n/x$ by $n_0\in\cqN$; its image is in the stated
kernel. Conversely any kernel representative $f$ has
$xf\in\cqN\cap x\cqE$. Finally $[n/x]=0$ exactly when
$n\in x\cqN$. This proves every asserted kernel and image.

The cokernel has only two possibilities. When every $n\in\cqN$
vanishes at $\lambda$, it is one-dimensional, with its exact
identification given by evaluation. Otherwise retain any
$h_\lambda\in\cqN$ with $h_\lambda(\lambda)\ne0$. Then
\begin{equation}\label{eq:cq-source-division}
 \begin{aligned}
 A_\lambda f(s)&=
  \frac{f(s)-f(\lambda)h_\lambda(s)/h_\lambda(\lambda)}{s-\lambda},\\
 f&=\frac{f(\lambda)}{h_\lambda(\lambda)}h_\lambda
                                      +(s-\lambda)A_\lambda f
 \end{aligned}
\end{equation}
proves surjectivity of $S-\lambda$, and its exact remaining kernel is
\begin{equation}\label{eq:cq-source-splitting}
 \ker(S-\lambda)\simeq
 \cqN/\bigl((s-\lambda)\cqN\oplus\cqC h_\lambda\bigr),
 \qquad n\longmapsto[A_\lambda n].
\end{equation}
The map is onto because a kernel class $[g]$ is the image of
$n=(s-\lambda)g\in\cqN$. Its kernel consists exactly of
$n=(s-\lambda)n_0+c h_\lambda$, $n_0\in\cqN$, by
\eqref{eq:cq-source-division}. The sum is direct by evaluation.
In this second case algebraic invertibility is therefore equivalent
to the exact identity
$\cqN=(s-\lambda)\cqN\oplus\cqC h_\lambda$.
On that locus the representative formula $[f]\mapsto[A_\lambda f]$
is independent of representatives, since $A_\lambda\cqN\subseteq
\cqN$; multiplication in \eqref{eq:cq-source-division} proves the
right inverse, and $A_\lambda((s-\lambda)f)=f$ proves the left
inverse. These are algebraic statements; continuity of this inverse
in $\tau$ has not been inserted.

The first derivative and the heat generator have exactly the same
frozen invariance question for the source subspace:
\begin{equation}\label{eq:cq-source-D-equivalence}
 D\cqN\subseteq\cqN
 \quad\Longleftrightarrow\quad D^2\cqN\subseteq\cqN
 \quad\Longleftrightarrow\quad -\tfrac14D^2\cqN\subseteq\cqN.
\end{equation}
Indeed the product rule is
$D^2(sn)-sD^2n=2Dn$. Since $sn\in\cqN$, second-derivative
invariance puts both terms on the left in $\cqN$, proving
first-derivative invariance. Iteration proves the reverse implication;
the retained nonzero scalar $-1/4$ proves the last equivalence.
If these equivalent inclusions hold, all $D^j\Xi$ belong to
$\cqN$. At each $\lambda$ at least one is nonzero: otherwise
the Taylor series would make $\Xi$ identically zero. Thus every
$S-\lambda$ is then algebraically surjective by
\eqref{eq:cq-source-division}. Its prescribed topological spectral
points would be represented entirely by the kernel
\eqref{eq:cq-source-splitting} or a noncontinuous algebraic inverse.
This computes the precise alternatives without assigning either
one to the source's stated spectral set.

There is also a full multiplicity consequence of algebraic derivative
descent. Writing $d[f]=[Df]$ on its exact descent locus, the product
rule gives $[d,S]=I$. For $0\ne v\in\ker(S-\lambda)$,
\begin{equation}\label{eq:cq-source-weylladder}
 (S-\lambda)d^kv=-k d^{k-1}v,
 \qquad (S-\lambda)^kd^kv=(-1)^k k!v\quad(k\geq0),
\end{equation}
with zero right side in the first identity at $k=0$.
To prove induction, commute the next $d$ using
$(S-\lambda)d=d(S-\lambda)-I$; the coefficient $-k$ becomes
$-(k+1)$. Iteration proves the second identity. A finite linear
dependence among $v,dv,\ldots,d^nv$, with largest nonzero
coefficient at index $j$, becomes that coefficient times
$(-1)^j j!v$ after applying $(S-\lambda)^j$, a contradiction.
Thus every such eigenvector generates an infinite independent
generalized eigenvector sequence. No finite source multiplicity
assertion has been assumed.

Finally even a single forward nonzero heat inclusion has a calculated
consequence. Fix $t\ne0$ and $c=t/2$, and define
\begin{equation}\label{eq:cq-source-heatcommutator}
 \begin{gathered}
 A_0=U_t,\qquad A_{k+1}=MA_k-A_kM,\qquad A_k=P_k(D)U_t,\\
 P_0(X)=1,\qquad P_{k+1}(X)=cX P_k(X)-P_k'(X).
 \end{gathered}
\end{equation}
The identity follows by induction from $[M,U_t]=cDU_t$ and
$[M,D^j]=-jD^{j-1}$: commuting $M$ through $P_k(D)U_t$
contributes $-P_k'(D)U_t+cP_k(D)DU_t$.
The degree of $P_k$ is exactly $k$, with leading coefficient
$c^k\ne0$. Under $U_t\cqN\subseteq\cqN$, induction with
$M\cqN\subseteq\cqN$ gives $A_k\cqN\subseteq\cqN$;
subtracting the lower-degree terms of $P_k$ then gives exactly
\begin{equation}\label{eq:cq-source-heatconsequence}
 D^kU_t\cqN\subseteq\cqN\quad(k\geq0).
\end{equation}
In particular every derivative of the nonzero $Z_t=U_t\Xi$ lies
in $\cqN$. At every point one derivative is nonzero by the entire
Taylor theorem, so every $S-\lambda$ is algebraically surjective
in this case too. This proof uses only the one forward inclusion.
For the stronger equality $U_t\cqN=\cqN$, it also proves
$D\cqN\subseteq\cqN$, by writing each $n\in\cqN$ as
$U_tm$ in \eqref{eq:cq-source-heatconsequence}. None of these
exact implications presumes the inclusion or equality for the actual
source closure. The explicit obstruction spaces above remain the
calculated source question.
